% =====================================================================
\chapter{Rendszerarchitektúra és helyi üzemeltetés}
\label{ch:architektura}
% =====================================================================

Ebben a fejezetben a \Mimir{} rendszer felépítését mutatjuk be abból a szempontból, amely a
kutatási kérdések szempontjából releváns: hogyan valósítja meg a rendszer a helyi futtatást (P4), a
zéró adatmegőrzést (P3) és a mérhetőséget. A felhasználói felület és az exportálók részleteit csak a
teljesség kedvéért érintjük.

\section{Tervezési célok}

A P1--P4 problémákból és az AI Act emberi felügyeletre vonatkozó elvárásából~\cite{aiact2024} öt
tervezési célt vezettünk le. \Az{\ref{tab:celok}}.~táblázat ezeket a közvetlen megközelítéssel veti össze.

\begin{table}[htbp]
\centering
\caption{Tervezési célok: a dokumentum csevegő LLM-nek való átadása és a \Mimir{} összevetése.}
\label{tab:celok}
\small
\begin{tabularx}{\textwidth}{@{}>{\raggedright\arraybackslash}p{2.6cm}>{\raggedright\arraybackslash}X>{\raggedright\arraybackslash}X@{}}
\toprule
Cél & Dokumentum $\to$ csevegő LLM & \Mimir{} \\
\midrule
G1 Ellenőrizhető & A kulcs nincs ellenőrizve, nincs forráshivatkozás & Minden kérdés darabokra hivatkozik; alátámasztottsági és vak megoldási ellenőrzés; az ellenőrzésen át nem jutott kérdés megjelölve \\
G2 Kontrollált & Darabszám, típus, nehézség csak kérés & A terv rögzíti a darabszámot, típusokat, Bloom-szinteket; lefedettségvezérelt fogalomválasztás \\
G3 Adatvédő & Az anyag a szolgáltatóhoz kerül & Helyi modell; a tartalom a feladat végén törlődik; kapcsolható kizárólag helyi mód \\
G4 Megvalósítható & Hosszú kontextus, nagy modell & Rövid hívások 8\,192--16\,384 tokenes kontextusban; 7B modell 8\,GB-on \\
G5 Felügyelt & A kimenet változtatás nélkül használható & Kötelező oktatói jóváhagyás szerkesztőben, a hivatkozott forrással; MI-közreműködés jelölése az exportban \\
\bottomrule
\end{tabularx}
\end{table}

\section{Mikroszolgáltatás-architektúra}

A rendszer lazán csatolt, konténerizált szolgáltatásokból áll, amelyeket egy Nginx-alapú API-átjáró
fog össze (\ref{fig:architektura}.~ábra). A szolgáltatások az északi mitológia tudással kapcsolatos
alakjairól kapták nevüket; feladatukat \az{\ref{tab:szolgaltatasok}}.~táblázat foglalja össze. Minden
szolgáltatás FastAPI-alapú Python-alkalmazás, egységes \code{/health} végponttal és verziózott
\code{/api/v1} felülettel, így a kiértékelő keretrendszer ugyanazokat a végpontokat hívja, mint az
éles frontend.

\begin{figure}[htbp]
\centering
\resizebox{\textwidth}{!}{%
\begin{tikzpicture}[
  svc/.style={draw, rounded corners=2pt, align=center, font=\footnotesize, minimum height=11mm, text width=25mm, fill=blue!6, draw=blue!50!black},
  ui/.style={svc, fill=white, draw=black},
  store/.style={draw, cylinder, shape border rotate=90, aspect=0.18, align=center, font=\footnotesize, minimum height=11mm, text width=19mm, fill=gray!10},
  ext/.style={svc, dashed, fill=orange!8, draw=orange!70!black},
  gpu/.style={svc, fill=green!8, draw=green!45!black, thick},
  arr/.style={-{Stealth[length=1.6mm]}, semithick},
  darr/.style={{Stealth[length=1.6mm]}-{Stealth[length=1.6mm]}, semithick}]
% sor 0
\node[ui]  (fe)    at (0,0)    {Webes felület\\(React, TypeScript)};
\node[ui]  (gw)    at (3.4,0)  {API-átjáró\\(Nginx, TLS)};
\node[svc] (bif)   at (6.8,0)  {\textbf{Bifrost}\\RAG-karmester,\\feladatkezelés};
\node[svc] (heim)  at (10.2,0) {\textbf{Heimdall}\\séma- és\\minőségszűrés};
\node[ext] (genai) at (12.6,2.0) {Egyetemi GenAI-\\szerver\\(opcionális)};
% sor +1
\node[svc] (well)  at (5.1,2.0) {\textbf{Wellspring}\\szövegkinyerés\\(PDF, DOCX, OCR)};
\node[svc] (rune)  at (8.5,2.0) {\textbf{RuneCarver}\\szemantikus\\darabolás (E5)};
% sor -1
\node[svc]   (skald) at (0,-2.0)   {\textbf{Skald}\\export: JSON,\\Moodle XML, PDF};
\node[svc]   (forge) at (3.4,-2.0) {\textbf{The Forge}\\feladatsor,\\audit napló};
\node[gpu]   (oll)   at (6.8,-2.0) {\textbf{Ollama}\\Qwen2.5-7B Q4\_K\_M\\(egyetlen GPU)};
\node[store] (qd)    at (10.2,-2.0) {Qdrant\\(memóriában)};
% sor -2
\node[store] (minio) at (0,-3.9)   {MinIO (S3)};
\node[store] (pg)    at (3.4,-3.9) {PostgreSQL};
% élek
\draw[arr]  (fe) -- (gw);
\draw[arr]  (gw) -- (bif);
\draw[darr] (bif) -- (well);
\draw[darr] (bif) -- (rune);
\draw[darr] (bif) -- (heim);
\draw[darr] (bif) -- (oll);
\draw[darr] (bif.south east) -- (qd.north west);
\draw[arr]  (gw.south west) -- (skald.north east);
\draw[arr]  (bif.south west) -- (forge.north east);
\draw[darr] (forge) -- (pg);
\draw[darr] (skald) -- (minio);
\draw[arr, dashed, orange!70!black] (bif.north east) -- node[above, sloped, font=\scriptsize, pos=0.55]{ha nincs LOCAL\_ONLY} (genai.west);
\begin{scope}[on background layer]
\node[draw=gray!60, dashed, rounded corners, fit=(well)(rune)(bif)(heim)(oll)(qd), inner sep=2.5mm,
  label={[font=\scriptsize, text=gray!70!black]below:feladat-hatókörű tartalom (zéró megőrzés)}] {};
\end{scope}
\end{tikzpicture}}
\caption{A \Mimir{} szolgáltatásai és adatfolyama. A szaggatott keretben lévő komponensek a
feltöltött tartalmat kizárólag a feladat élettartama alatt, memóriában tartják; a nyelvi modell
az egyetlen GPU-t használó összetevő.}
\label{fig:architektura}
\end{figure}

\begin{table}[htbp]
\centering
\caption{A \Mimir{} szolgáltatásai és kutatási szempontból releváns felelősségeik.}
\label{tab:szolgaltatasok}
\small
\begin{tabularx}{\textwidth}{@{}>{\raggedright\arraybackslash}p{2.4cm}>{\raggedright\arraybackslash}X@{}}
\toprule
Szolgáltatás & Felelősség \\
\midrule
Wellspring & Szövegkinyerés PDF-ből (PyMuPDF), DOCX-ből és szövegfájlból. Oldalanként olvashatósági arányt
számol (értelmes karakterek aránya a nem szóköz karakterek között); $0{,}85$ alatt az oldalt képként
dolgozza fel: helyi módban Tesseract OCR-rel, egyébként látás--nyelvi modellel. \\
RuneCarver & Mondatfelbontás, ablakozott mondatbeágyazás (multilingual-E5-base), szemantikus
darabolás méretfüggő küszöbbel (\ref{sec:darabolas}.~szakasz). \\
Bifrost & A generálás karmestere: aszinkron feladatkezelés lejárati idővel, indexelés a vektortárba,
visszakeresés, promptépítés, modellválasztás (helyi vagy szerver), audit metaadatok küldése. \\
Heimdall & A modellválaszok JSON-sémájának és koherenciájának szűrése; a szervermodell
elérhetetlensége esetén helyi tartalék. \\
The Forge & PostgreSQL-alapú feladatsor zárolásmentes kiosztással; audit napló kizárólag
metaadatokkal és SHA-256 ujjlenyomatokkal; megőrzési idő utáni automatikus törlés. \\
Skald & Exportálás JSON, Moodle XML (CDATA, MathJax) és nyomtatható PDF formátumba; az MI-közreműködés
és a modell nevének feltüntetése. \\
\bottomrule
\end{tabularx}
\end{table}

\subsection{Feladatkezelés és párhuzamosság}

A generálás percekig tarthat (\ref{ch:eredmenyek}.~fejezet), ezért a Bifrost aszinkron feladatmodellt
használ: a \code{POST /api/v1/generate} kérés azonnal egy feladatazonosítót ad vissza, az állapot
pedig a \code{GET /api/v1/status/\{job\_id\}} végponton kérdezhető le, amely a haladást és a becsült
hátralévő időt is közli. A feladatok egy memóriabeli tárban élnek, amely minden hozzáféréskor törli a
lejárati időn ($\mathrm{TTL}=3600$\,s) túli bejegyzéseket.

A CPU- és GPU-igényes háttérmunkát a The Forge osztja ki a munkásoknak egy PostgreSQL-táblán alapuló
sorral. A kiosztás a \code{SELECT \dots{} FOR UPDATE SKIP LOCKED} záradékra~\cite{postgres2024skiplocked}
épül: minden munkás egy tranzakcióban kiválasztja a legrégebbi, még nem zárolt, várakozó feladatot,
zárolja, állapotát \emph{feldolgozás alatt}-ra állítja, majd a befejezés után \emph{kész}-re. A
\code{SKIP LOCKED} miatt a párhuzamos munkások nem várnak egymás zárjaira, és egy feladatot pontosan egy
munkás kap meg, külön üzenetközvetítő (pl.\ RabbitMQ) nélkül. Mivel a GPU egyetlen modellt futtat, a
generálási feladatok párhuzamossága szándékosan egy; a párhuzamos munkások a CPU-igényes kinyerést és
darabolást gyorsítják.

\section{Adatvédelem: zéró megőrzés}
\label{sec:zero}

A G3 cél megvalósítása az adatminimalizálás és a tárolási korlátozás GDPR-elvein~\cite{gdpr2016}
alapul. \Az{\ref{tab:eletciklus}}.~táblázat a feldolgozás során keletkező összes adatfajta életciklusát
összegzi.

\begin{table}[htbp]
\centering
\caption{Adatfajták életciklusa a \Mimir{} rendszerben.}
\label{tab:eletciklus}
\small
\begin{tabularx}{\textwidth}{@{}>{\raggedright\arraybackslash}p{3.4cm}>{\raggedright\arraybackslash}p{3.0cm}>{\raggedright\arraybackslash}X@{}}
\toprule
Adat & Tárolás helye & Élettartam \\
\midrule
Feltöltött fájl bájtjai & szolgáltatás memóriája & a kinyerési kérés végéig \\
Kinyert szöveg, darabok & szolgáltatás memóriája & a feladat végéig \\
Beágyazások & Qdrant, memóriában & a következő feladat indexelése előtt a gyűjtemény törlődik \\
Promptok, kontextus & szolgáltatás memóriája & a modellhívás végéig \\
Generált vizsga (vázlat) & feladattár & $\mathrm{TTL}=60$ perc, majd automatikus törlés \\
Jóváhagyott, mentett vizsga & MinIO/SQLite & beállítható megőrzési idő, oktatói törlés \\
Audit napló & PostgreSQL & csak metaadatok és ujjlenyomatok; megőrzési idő után törlés \\
\bottomrule
\end{tabularx}
\end{table}

Az audit napló kulcsszerepet játszik a reprodukálhatóságban és az elszámoltathatóságban, ugyanakkor
nem tartalmazhat tananyagot. Ezért a napló egy bejegyzése a
\begin{equation}
a=\bigl(\textit{job\_id},\;\textit{modell},\;\textit{promptverzió},\;\textit{minőségi pontszám},\;
h(u),\;h(p),\;h(x),\;|u|,\;|x|\bigr)
\end{equation}
rekord, ahol $u$ a felhasználói kérés, $p$ a teljes prompt, $x$ a visszakeresett kontextus, $h$ a SHA-256
lenyomatfüggvény, $|\cdot|$ pedig a karakterhossz. A lenyomat alapján utólag igazolható, hogy két futás
ugyanazzal a kontextussal és prompttal dolgozott-e, a tartalom azonban nem állítható vissza belőle. A
rendszer induláskor a korábbi, még tartalmat is tároló naplósorok szöveges mezőit kiüríti, és óránként
törli a megőrzési időn túli sorokat.

A \code{LOCAL\_ONLY} kapcsoló minden külső modellhívást letilt: ilyenkor a látás--nyelvi modell helyett
Tesseract OCR, a szervermodell helyett a helyi Ollama fut, a modellválasztó végpont pedig csak helyi
modellt kínál. Kapcsoló nélkül a feltöltési képernyő megnevezi, hol történik a feldolgozás, és az oktató
választhatja a helyi modellt.

\section{Erőforrás-modell: a GPU-memória költségvetése}
\label{sec:vram}

A G4 cél kvantitatív kérdés: befér-e a generátor egy 8\,GB-os kártyára úgy, hogy mellette az asztali
környezet is fut? A dekódolás közbeni memóriaigényt három tag határozza meg:
\begin{equation}
M(\mathrm{ctx}) \;=\; W \;+\; \mathrm{ctx}\cdot \underbrace{2\,L\,H_{kv}\,d_h\,\beta}_{\textstyle m_{kv}} \;+\; M_0,
\label{eq:vram}
\end{equation}
ahol $W$ a kvantált súlyok mérete, $L$ a rétegek, $H_{kv}$ a kulcs--érték fejek száma, $d_h$ a fejdimenzió,
$\beta$ a gyorsítótár elemeinek mérete bájtban (fp16 esetén $\beta=2$), a $2$-es szorzó a kulcs és az érték
tárolását jelöli, $M_0$ pedig a futtatókörnyezet és a számítási pufferek többletköltsége. A
csoportosított lekérdezésű figyelem~\cite{ainslie2023gqa} lényege éppen az, hogy $H_{kv}$ a
lekérdezőfejek számánál jóval kisebb. A Qwen2.5-7B~\cite{qwen25} architektúrájára ($L=28$, $H_{kv}=4$,
$d_h=128$) a tokenenkénti gyorsítótár
\begin{equation}
m_{kv}=2\cdot 28\cdot 4\cdot 128\cdot 2\ \text{B}=57\,344\ \text{B}=56\ \text{KiB},
\end{equation}
így $16\,384$ tokenes kontextusnál a KV-gyorsítótár $56\cdot 16\,384\ \text{KiB}=896\ \text{MiB}$.
\Az{\ref{tab:budget}}.~táblázat a modell alapján becsült memóriaigényt mutatja több lehetséges helyi
generátorra.

\begin{table}[htbp]
\centering
\caption{Becsült GPU-memóriaigény 8\,192 tokenes kontextusnál, fp16 KV-gyorsítótárral és
$M_0\approx0{,}6$\,GB futtatókörnyezeti többlettel. A súlyok a GGUF-fájlok mérete, a gyorsítótár a
publikált architektúrákból számolt érték.}
\label{tab:budget}
\small
\begin{tabular}{@{}lcccccc@{}}
\toprule
Modell & Kvantálás & $L/H_{kv}$ & Súlyok & $m_{kv}$ & Összesen & Kártya \\
\midrule
Qwen2.5-3B  & Q4\_K\_M & 36/2 & 1,9\,GB & 36\,KiB  & 2,8\,GB  & 8\,GB \\
Qwen2.5-7B  & Q4\_K\_M & 28/4 & 4,7\,GB & 56\,KiB  & 5,8\,GB  & 8\,GB \\
Llama-3.1-8B~\cite{grattafiori2024llama3} & Q4\_K\_M & 32/8 & 4,9\,GB & 128\,KiB & 6,6\,GB & 8\,GB \\
Qwen2.5-7B  & Q8\_0    & 28/4 & 8,1\,GB & 56\,KiB  & 9,2\,GB  & 12\,GB \\
Qwen2.5-14B & Q4\_K\_M & 48/8 & 9,0\,GB & 192\,KiB & 11,2\,GB & 12\,GB \\
\bottomrule
\end{tabular}
\end{table}

A modell szerint a 3--8 milliárd paraméteres modellek 4 bites kvantálással elférnek egy 8\,GB-os
kártyán, a 14B modell vagy a 8 bites 7B modell viszont már 12\,GB-ot igényel. A becslést a pilot
vizsgálatban mért értékekkel validáltuk (\ref{tab:vram-meres}.~táblázat): minden kar első vizsgájánál,
amikor a modell még nem volt betöltve, az alapszint $1\,612$--$1\,931$\,MiB volt (asztali környezet), a
generálás alatti csúcs pedig $7\,022$--$7\,808$\,MiB. A modell és a futtatókörnyezet tehát
$5\,400$--$5\,950$\,MiB-ot foglalt $16\,384$ tokenes kontextus mellett, ami \az{\eqref{eq:vram}} modell alapján várt $4{,}7\,\text{GB}+0{,}9\,\text{GiB}+M_0$ értékkel összhangban van, és
a legrosszabb esetben is $\approx$380\,MiB tartalékot hagyott a $8\,192$\,MiB-os kártyán.

\begin{table}[htbp]
\centering
\caption{Mért GPU-memória a pilot vizsgálatban (MiB), karonként az első, hidegindítású vizsga alapján.
A különbség a modell és a futtatókörnyezet foglalása.}
\label{tab:vram-meres}
\small
\begin{tabular}{@{}lrrr@{}}
\toprule
Kar & Alapszint & Csúcs & Különbség \\
\midrule
\arm{E1}  & 1\,612 & 7\,025 & 5\,413 \\
\arm{E2}  & 1\,815 & 7\,761 & 5\,946 \\
\arm{E2f} & 1\,931 & 7\,621 & 5\,690 \\
\arm{E2h} & 1\,774 & 7\,175 & 5\,401 \\
\arm{E3}  & 1\,739 & 7\,134 & 5\,395 \\
\arm{E5a} & 1\,896 & 7\,369 & 5\,473 \\
\arm{E5b} & 1\,899 & 7\,326 & 5\,427 \\
\arm{E5c} & 1\,892 & 7\,292 & 5\,400 \\
\bottomrule
\end{tabular}
\end{table}

A mérés egy fontos mellékhatásra is rámutat: a teljes dokumentumot egy promptban kapó alapvonal
(\arm{B-doc-L}) számára a kontextus növelése nem ingyenes. $32\,768$ tokenes kontextusnál a
KV-gyorsítótár önmagában $1{,}75$\,GiB lenne, ami a mért csúcshoz adva már meghaladná a kártya
kapacitását; ekkor a futtatókörnyezet a modell egy részét a CPU-ra helyezi, ami nagyságrendi lassulással
jár. Ezért a helyi teljes dokumentumos alapvonal $16\,384$ tokenes kontextussal és $30\,000$ karakterre
csonkolt dokumentummal fut, és a levágott rész arányát rögzítjük -- éppen ez az a korlát, amelyet egy kis
helyi modell nem tud megkerülni, a tervezett csővezeték viszont igen, mert minden hívása rövid.
